\documentclass[../../../include/open-logic-section]{subfiles}

\begin{document}

\olfileid{sth}{choice}{banach}
\olsection{बानाख--टार्स्की विरोधाभास}

बूलोस यांनी प्रतिस्थापनाचे समर्थन करण्याचा प्रयत्न केला तसा, \emph{बाह्य} विचारांचा आधार घेऊन निवडीचे समर्थन करण्याचाही प्रयत्न करता येईल (पाहा \olref[replacement][extrinsic]{sec}). अखेर निवड स्वीकारल्याने अनेक इष्ट परिणाम मिळतात: प्रत्येक संचांकाची तुलना करता येते, प्रत्येक संचाचे सुक्रमण करता येते, संचांकांना विशिष्ट प्रकारच्या क्रमसूचक संख्या मानता येते, इत्यादी.

पण कधी कधी निवडीचे \emph{अनिष्ट} परिणाम आहेत, असा दावा केला जातो. मुख्यतः \cite{BanachTarski1924} यांच्या एका निष्कर्षामुळे हा दावा केला जातो.

\begin{thm}[बानाख--टार्स्की विरोधाभास ($\ZFC$ मध्ये)]
त्रिमितीय अवकाशातील धन त्रिज्येच्या कोणत्याही घनगोलाचे सांत संख्येने तुकडे करता येतात; घूर्णन आणि स्थानांतरणाने त्यांची पुनर्जुळणी करून त्या घनगोलाच्या दोन प्रती तयार करता येतात.
\end{thm}
\noindent
पहिल्या दृष्टीने हे थोडे थक्क करणारे आहे. या दोन घनगोलांचे घनफळ मूळ घनगोलाच्या \emph{दुप्पट} आहे. पण दृढ गती---घूर्णन आणि स्थानांतरण---घनफळ बदलत नाहीत. म्हणून बानाख--टार्स्कीमुळे जादूने नवे द्रव्य अस्तित्वात आणता येते, असे दिसते.

यापुढे आणखी आश्चर्य आहे.\footnote{पाहा \citet[Theorem 3.12]{Wagon2016}.} यासारखा युक्तिवाद दाखवतो की वाटाण्याचे सांत संख्येने तुकडे करून घूर्णन आणि स्थानांतरणाने त्यांची पुनर्जुळणी केल्यास बिग बेनच्या आकाराची आणि आकारमानाची वस्तू तयार करता येते.

पण केवळ $\ZF$ मध्ये यापैकी काहीही सिद्ध होत नाही.\footnote{तरी निवडीपेक्षा काटेकोरपणे दुर्बल तत्त्वांनी बानाख--टार्स्की सिद्ध करता येतो; पाहा \citet[303]{Wagon2016}.} म्हणून निवड नाकारावी की या “विरोधाभासासह” राहायला शिकावे, हा निर्णय घ्यावा लागतो.

या “विरोधाभासासह” राहायला शिकावे, असे आपण सुचवू. खरे तर हा फारसा विरोधाभास आहे, असे आम्हाला वाटत नाही. विशेषतः पुढील निष्कर्षांपेक्षा तो कमी किंवा अधिक विरोधाभासी का आहे, हे दिसत नाही:\footnote{\citet[276--7]{Potter2004}, \citet[16]{Weston2003} आणि \citet[31, 308--9]{Wagon2016} इतर उदाहरणांनी असेच मुद्दे मांडतात.}
\begin{enumerate}
	\item $(0,1)$ या अंतरालात जितके बिंदू आहेत, तितकेच $\Real$ मध्ये आहेत.
	\\\emph{पुरावा}: $\tan(\pi(r-\nicefrac{1}{2}))$ पाहा.
	\item रेषेत जितके बिंदू आहेत, तितकेच चौरसात आहेत.
	\\पाहा \olref[his][set][pathology]{sec} आणि \olref[his][set][cantorplane]{sec}.
	\item अवकाशभरण करणाऱ्या वक्ररेषा आहेत.
	\\पाहा \olref[his][set][pathology]{sec} आणि \olref[his][set][hilbertcurve]{sec}.
\end{enumerate}
या तिन्ही निष्कर्षांना निवडीची गरज नाही. आज आपण त्यांना गणिताचे आश्चर्यकारक, सुंदर भाग मानतो. बानाख--टार्स्की विरोधाभासाविषयीही तीच भूमिका घ्यावी, असे वाटू शकते.

येथे एक तांत्रिक निरीक्षण आवश्यक आहे; पण त्यासाठी शांतपणे विचार करणे पुरेसे आहे. दृढ गती घनफळ जतन करतात. त्यामुळे घनगोलाचे जे पाच\footnote{विरोधाभासाचे विधान “सांत संख्येने तुकडे” असे केले होते. खरे तर \citet{Robinson1947} यांनी \emph{पाच} तुकड्यांत (त्याहून कमी तुकड्यांत नाही) हे विघटन करता येते, असे सिद्ध केले. पुराव्यासाठी \citet[pp.~66--7]{Wagon2016} पाहा.} तुकडे केले जातात, ते सर्व \emph{मापनीय} असू शकत नाहीत. साधारणपणे, या स्वतंत्र तुकड्यांना घनफळ देण्यात अर्थ नाही. त्यांना चित्रात दाखवता न येणारे बिंदूंचे “अनंत विखुरणे” समजावे. अशा “अनंत विखुरलेल्या” संचांची कल्पना करणे विचित्र वाटू शकते. पण आधी उपलब्ध झालेल्या संचांची \emph{सर्व शक्य} संकलने तयार करावीत, या \stagesacc{} मध्ये व्यक्त झालेल्या निर्देशावरून त्यांचे अस्तित्व मिळत असल्यासारखे दिसते.

हे पटत नसेल, तर बानाख--टार्स्की विरोधाभास स्वीकारण्याच्या बाजूने शेवटचा एक बाह्य युक्तिवाद पाहा. त्यावरून गणितातील उत्कृष्ट विनोद लगेच मिळतो:
\begin{enumerate}
	\item[] \emph{प्रश्न}. “बानाख--टार्स्की”च्या अक्षरांची पुनर्रचना केल्यावर काय मिळते?
	\item[] \emph{उत्तर}. “बानाख--टार्स्की बानाख--टार्स्की”.
\end{enumerate}

\end{document}
